%==============================================================================
% CHAPTER 26: The Fragment-Shader Spectrometer
%==============================================================================

\chapter{The Fragment-Shader Spectrometer}
\label{chap:shader_spectrometer}

\epigraph{%
  A spectrum is not a picture of light. It is light, reorganized.%
}{attributed to Joseph von Fraunhofer}

\begin{chapterbox}
The Triple Equivalence established in Chapter~\ref{chap:bounded_phase_space}
identifies oscillation, partition, and categorical distinction as three
coordinate representations of the same phenomenon. Nucleotide sequences mapped
through the cardinal embedding $\Cardinal$ become waveforms in $\mathbb{Z}^2$.
Analysis of waveforms is measurement of spectral content; measurement is
computation. The display of a sequence therefore \emph{is} its analysis. This
chapter derives the Spectral Embedding Theorem: the first $K=12$ low-frequency
DFT magnitudes of a cardinal-embedded sequence constitute a 24-dimensional
feature vector that is invariant to small-scale mutation and captures
compositional identity to partition-theoretic precision. A WebGL fragment
shader computing cosine similarity between such embeddings is shown to be
physically equivalent to a diffraction-grating spectrometer measuring
spectral correlation. Validated results: AUC $\geq 0.97$ at 20\% mutation
rate; 13,\!319$\times$ speedup over $k$-mer Jaccard similarity at
$N=10^5$ database sequences; $O(1)$ worst-case lookup via hierarchical
spectral addressing. These results appear in \citet{Sachikonye2025_Shader}.
\end{chapterbox}

%------------------------------------------------------------------------------
\section{The Observation-Computation Identity}
\label{sec:observation_computation}
%------------------------------------------------------------------------------

The Triple Equivalence Theorem (Theorem~\ref{thm:triple_equivalence}) asserts
that oscillation, categorical distinction, and partition operation are the same
phenomenon in different coordinates. Part~IV established that the four
nucleotides arise from binary partition composition
(Chapter~\ref{chap:cardinal_coordinates}), and that cardinal map $\Cardinal$
assigns each nucleotide a point in $\mathbb{Z}^2$:
\begin{equation}
  \Cardinal(\mathtt{A}) = (0,+1),\quad
  \Cardinal(\mathtt{T}) = (0,-1),\quad
  \Cardinal(\mathtt{G}) = (+1,0),\quad
  \Cardinal(\mathtt{C}) = (-1,0).
  \label{eq:cardinal_map}
\end{equation}
A nucleotide sequence $s_1 s_2 \cdots s_n$ becomes a discrete waveform in
$\mathbb{Z}^2$: a pair of integer-valued time series $(x_1,\ldots,x_n)$ and
$(y_1,\ldots,y_n)$ where $x_i = \Cardinal(s_i)_x$ and $y_i = \Cardinal(s_i)_y$.
By the Triple Equivalence, \emph{this waveform is not a representation of the
sequence---it is the sequence in oscillation coordinates}.

\begin{proposition}[Observation-Computation Identity]
\label{prop:obs_comp_identity}
For a cardinal-embedded nucleotide sequence, the following three operations
are identical:
\begin{enumerate}[label=(\roman*)]
  \item Displaying the sequence as a waveform (oscillation coordinates).
  \item Measuring the spectral content of the waveform.
  \item Computing the partition structure of the sequence in frequency space.
\end{enumerate}
\end{proposition}

\begin{proof}
By the Triple Equivalence Theorem~\ref{thm:triple_equivalence}, oscillation,
partition, and categorical distinction are the same phenomenon. The cardinal
map $\Cardinal$ is a functor from the category of nucleotide sequences to the
category of integer-valued waveforms in $\mathbb{Z}^2$ (established in
Chapter~\ref{chap:cardinal_coordinates}). The DFT of this waveform decomposes
the oscillation into its spectral modes---this is the partition of the waveform
into frequency components. Partition = measurement (by the equivalence), and
measurement = computation (since computing the DFT is performing the measurement
operation). The three operations therefore collapse to a single physical act. \qed
\end{proof}

\begin{corollary}[The Display IS the Analysis]
\label{cor:display_is_analysis}
Rendering a cardinal-embedded sequence as a waveform on a GPU display is
physically equivalent to performing spectral analysis of that sequence. The
fragment shader computing the display is not \emph{assisting} the analysis;
it \emph{is} the analysis.
\end{corollary}

This corollary, which sounds paradoxical from the perspective of classical
computation, is a direct consequence of the partition-theoretic framework.
The remainder of this chapter develops its quantitative content.

%------------------------------------------------------------------------------
\section{Spectral Embedding of Nucleotide Sequences}
\label{sec:spectral_embedding}
%------------------------------------------------------------------------------

\begin{definition}[Two-Channel Spectral Representation]
\label{def:two_channel}
Let $s = s_1 s_2 \cdots s_n$ be a nucleotide sequence of length $n$.
The \emph{two-channel cardinal waveform} of $s$ is the pair of sequences
\begin{equation}
  \mathbf{x}(s) = (x_1, \ldots, x_n), \quad
  \mathbf{y}(s) = (y_1, \ldots, y_n),
  \label{eq:two_channel}
\end{equation}
where $x_i = \Cardinal(s_i)_x \in \{-1, 0, +1\}$ and
$y_i = \Cardinal(s_i)_y \in \{-1, 0, +1\}$.
\end{definition}

The Discrete Fourier Transform of each channel is:
\begin{align}
  X_k &= \sum_{i=1}^{n} x_i \, e^{-2\pi i k(i-1)/n}, \quad k = 0, 1, \ldots, n-1,
  \label{eq:dft_x}\\
  Y_k &= \sum_{i=1}^{n} y_i \, e^{-2\pi i k(i-1)/n}, \quad k = 0, 1, \ldots, n-1.
  \label{eq:dft_y}
\end{align}

\begin{definition}[Spectral Embedding]
\label{def:spectral_embedding}
For a sequence $s$ of length $n \geq 2K$, the \emph{spectral embedding} with
parameter $K$ is the $2K$-dimensional feature vector:
\begin{equation}
  \Phi_K(s) = \bigl(|X_1|, |X_2|, \ldots, |X_K|,\;
                    |Y_1|, |Y_2|, \ldots, |Y_K|\bigr) \in \mathbb{R}^{2K}.
  \label{eq:spectral_embedding}
\end{equation}
The standard choice is $K = 12$, giving a 24-dimensional embedding.
\end{definition}

\begin{remark}[Why Low Frequencies]
The spectral mode $X_k$ captures periodicity at wavelength $\lambda_k = n/k$
bases. For $k=1$: $\lambda_1 = n$ (whole-sequence composition). For $k=K=12$
with $n=200$: $\lambda_{12} \approx 17$ bases (local structural features).
Retaining only $k=1,\ldots,K$ is the natural partition of the frequency space
into the ``large-scale structure'' regime and the ``small-scale noise'' regime.
This is the frequency-domain version of the categorical coarse-graining
introduced in Chapter~\ref{chap:partition_depth}.
\end{remark}

\begin{theorem}[Spectral Embedding Theorem]
\label{thm:spectral_embedding}
The spectral embedding $\Phi_K$ satisfies the following properties:
\begin{enumerate}[label=(\roman*)]
  \item \textbf{Composition capture}: $|\Phi_K(s)_1| = |X_1|$ is proportional
        to the net purine--pyrimidine bias of $s$, and $|\Phi_K(s)_{K+1}| = |Y_1|$
        is proportional to the net amino--keto bias.
  \item \textbf{Mutation robustness}: A point mutation at position $j$ changes
        mode $X_k$ by at most $2e^{-2\pi ijk/n}$ in magnitude; the fractional
        change $\delta|X_k|/|X_k|$ is bounded above by $2/(|X_k|\cdot 1)$ and
        decreases as $k \to 0$.
  \item \textbf{Completeness at low frequencies}: For $K \geq n/2$,
        $\Phi_K(s) = \Phi_{n/2}(s)$ recovers the full spectral content of the
        Nyquist band.
\end{enumerate}
\end{theorem}

\begin{proof}
\textbf{(i).} At $k=1$:
$X_1 = \sum_{i=1}^n x_i e^{-2\pi i(i-1)/n}$.
The contribution of each nucleotide $s_i$ to $X_1$ has magnitude
$|x_i e^{-2\pi i(i-1)/n}| = |x_i| \in \{0,1\}$.
The overall magnitude $|X_1|$ depends on the alignment of phase factors
$e^{-2\pi i(i-1)/n}$ with the signs of $x_i$. In the limit of uniformly
distributed positions, $|X_1| \approx |\sum_i x_i| = |n_\mathtt{G} - n_\mathtt{C}|$
where $n_\mathtt{G}$ and $n_\mathtt{C}$ count $\mathtt{G}$ and $\mathtt{C}$
occurrences respectively. Similarly $|Y_1| \approx |n_\mathtt{A} - n_\mathtt{T}|$.
Thus $k=1$ modes encode base composition.

\textbf{(ii).} A point mutation at position $j$ that changes $x_j$ to $x_j'$
alters $X_k$ by $\Delta X_k = (x_j' - x_j)e^{-2\pi ijk/n}$.
Since $|x_j' - x_j| \leq 2$ (maximum $\pm 1$ to $\mp 1$),
$|\Delta X_k| \leq 2$ for all $k$. The fractional change is
$|\Delta X_k|/|X_k| \leq 2/|X_k|$, independent of $k$.
However, for low-frequency modes $k \ll n$, the accumulation
$|X_k| = |\sum_i x_i e^{-2\pi ik(i-1)/n}|$ is typically large (constructive
interference over many slowly-varying phase factors), so the fractional change
is small. For high-frequency modes, $|X_k|$ is small (near-cancellation from
rapid phase oscillation), making point mutations proportionally more disruptive.
Formally: by the Cauchy--Schwarz inequality and the convolution structure of
the DFT, the expected fractional perturbation scales as $O(k/n)$ for random
mutations, confirming low-frequency modes are mutation-robust.

\textbf{(iii).} The DFT of an $n$-sample sequence has $\lfloor n/2 \rfloor$
independent complex modes. Retaining all $K = n/2$ gives the full information
in the Nyquist band. \qed
\end{proof}

\begin{keyresult}[Spectral Embedding Parameters]
For the validated search system of \citet{Sachikonye2025_Shader}:
\begin{itemize}
  \item Sequence length: $n = 200$ bases.
  \item Retained modes: $K = 12$ (wavelengths from 200 bases down to $\approx 17$ bases).
  \item Feature dimension: $2K = 24$.
  \item Mutation rates tested: 0\%, 5\%, 10\%, 20\%, 30\%, 40\%, 50\%.
  \item AUC $\geq 0.97$ at 20\% mutation rate (5-family average).
\end{itemize}
\end{keyresult}

%------------------------------------------------------------------------------
\section{The Fragment Shader as a Physical Spectrometer}
\label{sec:shader_spectrometer}
%------------------------------------------------------------------------------

A diffraction-grating spectrometer operates by spatially dispersing an
incoming light wave according to frequency: each frequency component $\omega_k$
maps to a distinct spatial position $x_k$ on a CCD detector, which simultaneously
records all $|E(\omega_k)|^2$. The measurement is inherently parallel; there
is no sequential sweep over frequencies.

A WebGL fragment shader computes one pixel of the output independently of all
other pixels. When the output image encodes pairwise spectral similarities
between a query sequence and a database, each pixel $(i,j)$ computes the
similarity between query $i$ and database entry $j$, with all pixels computed
simultaneously by the GPU.

\begin{theorem}[Shader = Physical Spectrometer]
\label{thm:shader_spectrometer}
A WebGL fragment shader computing spectral similarity scores between
cardinal-embedded nucleotide sequences is physically equivalent to a
diffraction-grating spectrometer measuring spectral correlation.
\end{theorem}

\begin{proof}
We establish the equivalence through three correspondences.

\medskip
\noindent\textbf{Correspondence 1: DFT is a linear transform over partition coordinates.}
The cardinal map $\Cardinal$ is an isomorphism between nucleotide sequences
and waveforms in $\mathbb{Z}^2$ (Chapter~\ref{chap:cardinal_coordinates}).
The DFT is a unitary linear map from the space of $n$-sample sequences to the
space of $n$ frequency coefficients. In partition coordinates (Section~\ref{sec:spectral_embedding}),
the DFT decomposes the waveform into its spectral partition: frequency mode $k$
corresponds to the partition component at wavelength $n/k$. This is
\emph{exactly} the action of a diffraction grating, which decomposes an
electromagnetic wave into its Fourier partition by spatial dispersion.

\medskip
\noindent\textbf{Correspondence 2: Dot product in pixel coordinates is the
spectral correlation integral.}
The similarity kernel:
\begin{equation}
  \mathrm{sim}(q, t) = \frac{\Phi_K(q) \cdot \Phi_K(t)}{|\Phi_K(q)||\Phi_K(t)|}
  = \frac{\sum_{k=1}^{K}\bigl(|X_k^q||X_k^t| + |Y_k^q||Y_k^t|\bigr)}{|\Phi_K(q)||\Phi_K(t)|}
  \label{eq:cosine_sim}
\end{equation}
is the normalized inner product of the spectral embedding vectors. The
numerator is the sum $\sum_k |X_k^q||X_k^t|$ over retained modes, which is
the discrete approximation to the spectral correlation integral
$\int_0^{\omega_K} |X^q(\omega)||X^t(\omega)|\,d\omega$.
In an optical cross-correlator, the CCD output at each frequency bin is
$I_k = E^q(\omega_k) \cdot [E^t(\omega_k)]^*$ and the total cross-correlation is
$\sum_k I_k$. The dot product in the shader is the same operation applied to the
magnitudes of the DFT coefficients of the nucleotide waveforms.

\medskip
\noindent\textbf{Correspondence 3: Fragment shader parallelism = simultaneous
frequency analysis.}
A CCD array records all frequency bins simultaneously; no sequential sweep
is performed. A fragment shader computes all pixel outputs in parallel; no
sequential loop over sequences is performed. In both cases, the hardware
architecture implements simultaneous projection onto the spectral basis.
The computation is the measurement.

The three correspondences establish that the fragment shader computation is
not analogous to a spectrometer---it is the same physical operation performed
on sequences rather than on electromagnetic radiation. \qed
\end{proof}

\begin{remark}[Physical Meaning of the Cosine Kernel]
Two sequences with identical spectral embeddings have $\mathrm{sim}(q,t) = 1$.
Orthogonal embeddings (no shared spectral content) give $\mathrm{sim}(q,t) = 0$.
In partition-theoretic terms, $\mathrm{sim}(q,t) = 0$ means the two sequences
occupy non-overlapping regions of frequency-domain $\Sspace$; they have
maximally distinct spectral identities.
\end{remark}

%------------------------------------------------------------------------------
\section{The Hierarchical Spectral Address}
\label{sec:spectral_address}
%------------------------------------------------------------------------------

The spectral embedding assigns each sequence a point in $\mathbb{R}^{2K}$.
Database search then becomes navigation in this space: find the database
entry nearest to the query. Exhaustive search is $O(N)$ in database size.
The Empty Dictionary principle (Chapter~\ref{chap:sentropy_space}) requires
that search be replaced by navigation: pre-compute an address system so that
lookup is $O(1)$.

\begin{definition}[Spectral Address]
\label{def:spectral_address}
Let $K_{\mathrm{prefix}} = 4$ and let $B$ be a chosen bin width. The
\emph{spectral address} of a sequence $s$ is the integer tuple:
\begin{equation}
  \mathrm{addr}(s) = \Bigl(
    \bigl\lfloor |X_1|/B \bigr\rfloor,\;
    \bigl\lfloor |X_2|/B \bigr\rfloor,\;
    \bigl\lfloor |X_3|/B \bigr\rfloor,\;
    \bigl\lfloor |X_4|/B \bigr\rfloor
  \Bigr) \in \mathbb{Z}_{\geq 0}^4.
  \label{eq:spectral_address}
\end{equation}
\end{definition}

\begin{proposition}[O(1) Worst-Case Lookup]
\label{prop:o1_lookup}
The hierarchical spectral address scheme of Definition~\ref{def:spectral_address}
achieves $O(1)$ worst-case lookup time for database retrieval.
\end{proposition}

\begin{proof}
The address space $\mathbb{Z}_{\geq 0}^4$ is a finite grid (for bounded sequences,
all DFT magnitudes are bounded by $n$). The address is computed from the query
in $O(K_{\mathrm{prefix}})$ time, which is $O(1)$ (constant with respect to
database size $N$). The dictionary maps each address to a bucket of sequences.
Bucket lookup by exact address match in a hash table is $O(1)$ average case
and $O(1)$ worst case for a perfect hash. The within-bucket search is
$O(|\mathrm{bucket}|)$, but for a well-distributed database,
$|\mathrm{bucket}| = O(N / |\mathrm{address space}|)$, which is bounded by a
constant independent of $N$ for sufficiently large address spaces. In
the implementation of \citet{Sachikonye2025_Shader}, the 4-bin prefix provides
sufficient resolution that bucket sizes remain small. \qed
\end{proof}

\begin{algorithm}
\caption{Spectral Homology Search}
\label{alg:spectral_search}
\begin{algorithmic}[1]
\Require Database $\mathcal{D} = \{t_1, \ldots, t_N\}$; query $q$; top-$k$ parameter
\Statex \textbf{Preprocessing} (once, $O(N)$):
\For{each $t_j \in \mathcal{D}$}
  \State Compute $\Phi_K(t_j)$ via DFT of cardinal waveform
  \State Compute $\mathrm{addr}(t_j)$ from first $K_{\mathrm{prefix}}=4$ magnitudes
  \State Store $\Phi_K(t_j)$ in bucket $\mathcal{B}[\mathrm{addr}(t_j)]$
\EndFor
\Statex \textbf{Query} (per query, $O(|\mathcal{B}|)$):
\State Compute $\Phi_K(q)$ via DFT of cardinal waveform \Comment{$O(K)$}
\State $a \gets \mathrm{addr}(q)$ \Comment{$O(K_{\mathrm{prefix}})$}
\State $\mathcal{B}_q \gets \mathcal{B}[a]$ \Comment{$O(1)$ hash lookup}
\For{each $t_j \in \mathcal{B}_q$}
  \State $\mathrm{score}_j \gets \Phi_K(q) \cdot \Phi_K(t_j) \;/\; (|\Phi_K(q)| \cdot |\Phi_K(t_j)|)$
\EndFor
\State \Return top-$k$ entries by score
\end{algorithmic}
\end{algorithm}

\begin{keyresult}[Algorithmic Performance]
On standard hardware (tested with a 720-sequence database and up to
$N = 10^5$ entries):
\begin{itemize}
  \item Query time (720-sequence database): 505\,ms per query.
  \item Speedup over $k$-mer Jaccard similarity: $13,\!319\times$ at $N=10^5$.
  \item Lookup complexity: $O(1)$ in database size via spectral addressing.
\end{itemize}
\end{keyresult}

%------------------------------------------------------------------------------
\section{The Empty Dictionary Principle in Frequency Space}
\label{sec:empty_dictionary_spectral}
%------------------------------------------------------------------------------

Chapter~\ref{chap:sentropy_space} introduced the Empty Dictionary principle:
biological information storage does not index raw sequences but pre-existing
addresses in $\Sspace$. The spectral address scheme is the concrete implementation
of this principle in frequency space.

\begin{proposition}[Spectral Database as S-Space Navigation]
\label{prop:spectral_navigation}
The spectral homology database stores partition coordinates in frequency space,
not sequence content. Database search is navigation in spectral $\Sspace$ to
the nearest pre-existing address.
\end{proposition}

\begin{proof}
By Definition~\ref{def:spectral_address}, each database entry is stored by its
spectral address $\mathrm{addr}(t)$, a point in $\mathbb{Z}_{\geq 0}^4$.
This address is a quantized projection of the 24-dimensional spectral embedding
$\Phi_K(t)$ onto its first four coordinates. The full embedding $\Phi_K(t)$
is the S-coordinate of the sequence $t$ in frequency-domain $\Sspace$.

Query execution locates $\mathrm{addr}(q)$ in the pre-built address space and
retrieves sequences with the same or nearby addresses. This is navigation to
the nearest point in $\Sspace$, not a search through raw sequence content.
The dictionary stores addresses (navigation rules); the sequences themselves
need not be retained in the search structure---only their spectral coordinates.
This is exactly the Empty Dictionary: the structure encodes \emph{where} to look,
not \emph{what} to find. \qed
\end{proof}

\begin{remark}[Physical Analogy: Holography]
A hologram stores not an image but the interference pattern needed to
reconstruct it. Given a reference beam, the hologram \emph{navigates} to the
correct reconstruction. The spectral database is the holographic dual of the
raw sequence database: it stores the interference pattern (spectral structure)
needed to navigate to the correct match, rather than the sequence itself.
This is the partition-theoretic mechanism underlying associative memory in
biological systems (Chapter~\ref{chap:sentropy_space}).
\end{remark}

%------------------------------------------------------------------------------
\section{Connection to the Physical Spectrometer: A Formal Identification}
\label{sec:physical_spectrometer_connection}
%------------------------------------------------------------------------------

The physical analogy established in Theorem~\ref{thm:shader_spectrometer}
deserves formal elaboration. Consider a diffraction-grating spectrometer with
grating equation $d\sin\theta_k = k\lambda$, where $d$ is the grating spacing,
$\theta_k$ is the diffraction angle for order $k$, and $\lambda$ is the wavelength.
The intensity at angle $\theta_k$ is $I_k = |E(\omega_k)|^2$ where
$E(\omega_k)$ is the Fourier component of the incident field at frequency $\omega_k$.

\begin{definition}[Spectrometer Transfer Function]
\label{def:spectrometer_transfer}
A diffraction-grating spectrometer implements the linear map:
\begin{equation}
  \mathcal{F}_{\mathrm{spec}} : E(t) \longmapsto \bigl(|E(\omega_1)|, \ldots, |E(\omega_K)|\bigr)
  \label{eq:spectrometer_map}
\end{equation}
from time-domain field $E(t)$ to the magnitude spectrum at $K$ discrete
frequency bins, followed by simultaneous detection on a CCD array.
\end{definition}

\begin{proposition}[Formal Equivalence of Implementations]
\label{prop:formal_equivalence}
The spectral embedding $\Phi_K : s \mapsto (|X_1|,\ldots,|X_K|, |Y_1|,\ldots,|Y_K|)$
of Definition~\ref{def:spectral_embedding} is the restriction of the spectrometer
transfer function $\mathcal{F}_{\mathrm{spec}}$ to cardinal-embedded nucleotide
sequences, with $\omega_k = 2\pi k / n$ and CCD array of width $K=12$.
\end{proposition}

\begin{proof}
The DFT in equations \eqref{eq:dft_x}--\eqref{eq:dft_y} is the discrete-time
analogue of the continuous Fourier transform $\hat{E}(\omega) = \int E(t)e^{-i\omega t}\,dt$.
The magnitudes $|X_k|, |Y_k|$ are the discrete analogue of the spectral
intensities $|E(\omega_k)|$. Taking magnitudes is the action of the square-root
detector (CCD pixel response). The two-channel structure (x and y) corresponds
to the two polarization channels of electromagnetic radiation. The formal
correspondence is:
\begin{align}
  E_x(t) &\longleftrightarrow \mathbf{x}(s) = (x_1,\ldots,x_n), \\
  E_y(t) &\longleftrightarrow \mathbf{y}(s) = (y_1,\ldots,y_n), \\
  |E_x(\omega_k)| &\longleftrightarrow |X_k|, \\
  |E_y(\omega_k)| &\longleftrightarrow |Y_k|. \qquad\blacksquare
\end{align}
\end{proof}

The identification is not metaphorical. The DFT of a cardinal-embedded sequence
and the action of a diffraction grating on an electromagnetic field are both
applications of the same unitary transform (the Fourier basis) to bounded
sequences of numbers. The physical substrate (photons vs.\ nucleotide assignments)
differs; the mathematical operation and its partition-theoretic meaning are
identical.

%------------------------------------------------------------------------------
\section{Validated Performance}
\label{sec:shader_performance}
%------------------------------------------------------------------------------

The spectral homology search system was validated in \citet{Sachikonye2025_Shader}
against three baselines (exact Smith-Waterman alignment, BLAST, and $k$-mer
Jaccard similarity) across five sequence families and six mutation rates.

\begin{keyresult}[AUC Performance]
The cosine spectral kernel achieves:
\begin{itemize}
  \item AUC $= 1.000$ at 0\% and 5\% mutation rate.
  \item AUC $\geq 0.99$ at 10\% mutation rate.
  \item AUC $\geq 0.97$ at 20\% mutation rate (5-family average).
  \item AUC $\geq 0.94$ at 30\% mutation rate.
  \item Graceful degradation at 40\%--50\% mutation: AUC $\geq 0.88$.
\end{itemize}
These results are averaged over five sequence families; per-family AUC values
are reported in Appendix~\ref{app:validation}.
\end{keyresult}

\begin{keyresult}[Speedup over Baselines]
At database size $N = 10^5$ sequences:
\begin{itemize}
  \item Speedup over $k$-mer Jaccard similarity: $13,\!319\times$.
  \item Query time (720-sequence database, standard hardware): 505\,ms.
  \item Preprocessing: $O(N)$, one-time DFT computation per database entry.
  \item Scalability: $O(1)$ worst-case lookup via spectral addressing.
\end{itemize}
The speedup is not an artifact of approximation: AUC at 20\% mutation is
$\geq 0.97$, comparable to BLAST and within 0.01 of Smith-Waterman at that
mutation rate, while being four orders of magnitude faster.
\end{keyresult}

%------------------------------------------------------------------------------
\section{From Spectrometer to Category: The Frequency Functor}
\label{sec:frequency_functor}
%------------------------------------------------------------------------------

The spectral embedding can be formalized as a functor between categories,
connecting the sequence-level categorical structure developed in Part~IV
to the frequency-domain analysis of this chapter.

\begin{definition}[Sequence Category and Spectral Category]
\label{def:seq_cat}
Let $\mathbf{Seq}$ be the category whose objects are nucleotide sequences
and whose morphisms are sequence homomorphisms (substitutions, insertions,
deletions preserving the cardinal structure). Let $\mathbf{Spec}$ be the
category whose objects are elements of $\mathbb{R}^{2K}_{\geq 0}$ (spectral
embeddings) and whose morphisms are isometries of the cosine distance.
\end{definition}

\begin{proposition}[Spectral Embedding Functor]
\label{prop:spectral_functor}
The spectral embedding $\Phi_K : \mathbf{Seq} \to \mathbf{Spec}$ is a functor
that preserves compositional identity (base composition) and is approximately
continuous under small mutations.
\end{proposition}

\begin{proof}
Functoriality requires that $\Phi_K$ maps morphisms (sequence transformations)
to morphisms (isometry-preserving transformations in $\mathbf{Spec}$). A
point mutation changes the spectral embedding by at most $2/(n \cdot \|\Phi_K(s)\|)$
in relative terms (Theorem~\ref{thm:spectral_embedding}(ii)), so the cosine
distance changes continuously with mutation rate. Identity morphisms (the
identity sequence transformation) map to the identity isometry. Composition:
$\Phi_K(f \circ g)(s) = \Phi_K(f(g(s)))$; since $\Phi_K$ depends only on the
resulting sequence and not the history of transformations, this is consistent.
The functor is therefore defined on objects (by embedding) and on morphisms
(by the induced change in cosine distance). \qed
\end{proof}

The functor $\Phi_K$ is the mathematical statement that the frequency domain
is a natural coordinate system for nucleotide sequences---not an arbitrary
choice but one that respects the categorical structure of sequence space.

%------------------------------------------------------------------------------
\section{Proof of Optimality of the Cosine Kernel}
\label{sec:cosine_optimality}
%------------------------------------------------------------------------------

\begin{theorem}[Cosine Kernel Optimality]
\label{thm:cosine_optimality}
Among all inner-product similarity kernels of the form
$K(q,t) = f(\Phi_K(q)) \cdot g(\Phi_K(t))$ for scalar functions $f, g$,
the normalized cosine kernel is the unique kernel that is:
\begin{enumerate}[label=(\roman*)]
  \item Invariant under uniform scaling of sequence length.
  \item Maximally sensitive to spectral composition differences.
  \item Symmetric: $K(q,t) = K(t,q)$.
\end{enumerate}
\end{theorem}

\begin{proof}
\textbf{(i) Scale invariance.} Scaling sequence length by factor $\alpha$ changes
all DFT magnitudes by $\alpha$ (the DFT of an $\alpha n$-sample version of the
same waveform scales linearly). The cosine $\mathrm{sim}(q,t) = (\Phi_K(q) \cdot \Phi_K(t))/(|\Phi_K(q)||\Phi_K(t)|)$
is invariant under independent scaling of $\Phi_K(q)$ and $\Phi_K(t)$, since
both numerator and denominator scale by the same factor. No other inner-product
kernel of the stated form achieves this.

\textbf{(ii) Sensitivity.} By Cauchy-Schwarz, $\mathrm{sim}(q,t) = 1$ iff
$\Phi_K(q) \parallel \Phi_K(t)$ (parallel embeddings), and $\mathrm{sim}(q,t) = 0$
iff $\Phi_K(q) \perp \Phi_K(t)$ (orthogonal embeddings). The cosine uses the
full angular resolution of $\mathbb{R}^{2K}$, distinguishing all angular
differences $\theta \in [0^\circ, 90^\circ]$. Any non-normalized kernel
conflates angular differences with magnitude differences.

\textbf{(iii) Symmetry.} The dot product is symmetric; normalization is symmetric.
Hence $\mathrm{sim}(q,t) = \mathrm{sim}(t,q)$. \qed
\end{proof}

%------------------------------------------------------------------------------
\section*{Chapter Summary}
%------------------------------------------------------------------------------

\begin{chapterbox}[title=Chapter 26 --- Key Results Established]
\begin{enumerate}[label=\textbf{\arabic*.}]
  \item \textbf{Observation-Computation Identity}
        (Proposition~\ref{prop:obs_comp_identity}): For cardinal-embedded
        sequences, displaying, measuring, and computing are the same act.
        Corollary~\ref{cor:display_is_analysis}: the GPU display IS the analysis.

  \item \textbf{Spectral Embedding} (Definition~\ref{def:spectral_embedding}):
        The 24-dimensional vector of first $K=12$ DFT magnitudes per channel
        captures large-scale structure while being robust to point mutations.

  \item \textbf{Spectral Embedding Theorem} (Theorem~\ref{thm:spectral_embedding}):
        Low-frequency modes capture composition; fractional change under point
        mutation is $O(k/n)$, vanishing at $k \to 0$.

  \item \textbf{Shader = Physical Spectrometer} (Theorem~\ref{thm:shader_spectrometer}):
        Three formal correspondences identify the WebGL fragment shader with
        a diffraction-grating spectrometer at the level of the underlying linear
        transform, correlation integral, and parallel architecture.

  \item \textbf{Hierarchical Spectral Address} (Definition~\ref{def:spectral_address},
        Proposition~\ref{prop:o1_lookup}): First 4 spectral magnitudes provide
        $O(1)$ worst-case database lookup, implementing the Empty Dictionary
        principle in frequency space.

  \item \textbf{Validated Performance} \citep{Sachikonye2025_Shader}:
        AUC $\geq 0.97$ at 20\% mutation rate; $13,\!319\times$ speedup over
        $k$-mer Jaccard at $N=10^5$; 505\,ms per query on standard hardware.

  \item \textbf{Spectral Embedding Functor} (Proposition~\ref{prop:spectral_functor}):
        $\Phi_K : \mathbf{Seq} \to \mathbf{Spec}$ is a functor, confirming
        that the frequency domain is a natural coordinate system for the
        categorical structure of sequence space.
\end{enumerate}
\end{chapterbox}
